% TOP_PROBLEM: {"id": 3075, "title": "Average diameter of a bounded cell of a simple arrangement", "category_id": 6, "category": {"id": 6, "name": "geometry", "display_name": "Geometry", "description": "Euclidean and non-Euclidean geometry, geometric structures.", "slug": "geometry", "order_index": 6, "created_at": "2026-07-31T15:26:25.671Z"}, "status": "open", "problem_number": "OPG-605", "set_id": 12}
% FIRST_POSTED: 2026-10-01
% ATTEMPT_STATUS: unresolved
% UnsolvedMath ID 3075; source catalogue code OPG-605
% !TeX program = lualatex
% GPT-6 Astra Ultra research attempt; independent partial work, 2026-10-01.
\documentclass[11pt,a4paper]{article}
\usepackage[margin=25mm]{geometry}
\usepackage{fontspec}
\setmainfont{lmroman10-regular.otf}[BoldFont=lmroman10-bold.otf,
 ItalicFont=lmroman10-italic.otf,BoldItalicFont=lmroman10-bolditalic.otf]
\usepackage{amsmath,amssymb,mathtools}
\usepackage{unicode-math}
\setmathfont{latinmodern-math.otf}
\AtBeginDocument{\let\setminus\smallsetminus}
\usepackage{xurl,hyperref}
\hypersetup{colorlinks=true,urlcolor=blue}
\setcounter{secnumdepth}{0}
\setlength{\parindent}{0pt}
\setlength{\parskip}{5pt}
\setlength{\emergencystretch}{3em}
\begin{document}
\section*{Average diameter of a bounded cell of a simple arrangement}\label{3075}
{\small UnsolvedMath ID 3075 · OPG-605 · Geometry · OpenGarden}
\subsection{Definitions and mathematical statement}
Fix integers $d\ge1$ and $n\ge d+1$. A simple affine hyperplane
arrangement in $\mathbb R^d$ is a family $\mathcal A$ of $n$ hyperplanes
in general position: every $j\le d$ distinct hyperplanes meet in an
affine subspace of dimension $d-j$, and no $d+1$ have a common point.
The components of their complement are open polyhedral regions.
The closures of the bounded regions are full-dimensional convex polytopes,
called the bounded cells.

For such a cell $P$, let $\delta(P)$ be the diameter of its edge graph:
the maximum, over pairs of vertices, of the smallest number of polytope
edges in a path joining them. This is a combinatorial graph distance,
not a Euclidean distance between points. List all bounded cells as
$P_1,\ldots,P_I$ and define
\[
 \overline\delta(\mathcal A)=\frac1I\sum_{j=1}^I\delta(P_j).
\]
The general-position assumptions give $I=\binom{n-1}{d}>0$, so the
average is defined and weights every bounded cell equally.

The conjecture states that
$\overline\delta(\mathcal A)\le d$ for every such arrangement, independently
of $n$. Individual cells are allowed to have diameter greater than $d$;
the requirement bounds only the mean over all bounded cells.
Unbounded regions are omitted, and neither cell volume nor number of
vertices changes its weight in the average. Establishing a bound for just
one selected polytope therefore does not settle the arrangement statement.
\subsection{Short English statement}
Is the average graph diameter of all bounded cells of a simple hyperplane arrangement at most the ambient dimension?
\subsection{Research attempt: planar incidence accounting and its dimensional limit}
\paragraph{Scope and literature check.}
GPT-6 Astra Ultra, 1 October 2026. Diameters are graph distances and
the average gives each bounded cell equal weight. Deza--Xie [S3]
prove the conjecture in dimension two and determine its exact extremal
value there; they also give higher-dimensional estimates and constructions.
Thus the planar bound below is a reconstruction of a weaker elementary
estimate, not an improvement or a newly solved case. The Garden URL
[S2] did not load in this check, so the research paper is used for
the literature qualification.

\paragraph{1. An elementary bound tending to two in the plane.}
Let $d=2$ and $n\ge3$. Each of the $n$ lines has $n-1$ distinct
intersection points, because no two lines are parallel and no three
are concurrent. These divide the line into two rays and exactly
$n-2$ bounded elementary segments. Thus the arrangement has
\[
 E_b=n(n-2)
\]
bounded edges. Every side of a bounded polygonal cell is one of these
edges, and each edge borders at most two bounded cells. If the
bounded cells have $m_1,\ldots,m_I$ sides, then
$\sum_jm_j\le2E_b$.

The edge graph of a convex $m$-gon is a cycle, of diameter
$\lfloor m/2\rfloor$. To check this formula, the two routes between
any pair of vertices have lengths summing to $m$, so the shorter is
at most $\lfloor m/2\rfloor$; a pair separated by that many cycle
edges attains it. Using $I=(n-1)(n-2)/2$, we obtain
\[
 \overline\delta\le\frac1{2I}\sum_jm_j
 \le\frac{n(n-2)}I=\frac{2n}{n-1}
 =2+\frac2{n-1}.
\]
This controls the whole average; no bound on the largest individual
polygon diameter was assumed. It remains too large by $2/(n-1)$
to prove the exact desired inequality.

\paragraph{2. Isolate the precise missing incidence saving.}
For a bounded edge $e$, let $b(e)\in\{0,1,2\}$ be the number of
bounded cells incident to it. Define
$D=\sum_e(2-b(e))=2E_b-\sum_jm_j$, a nonnegative deficit.
Let $O$ be the number of bounded cells with an odd number of sides.
The cycle formula gives the exact identity
\[
 \overline\delta
 =\frac{\sum_jm_j-O}{2I}
 =\frac{2n(n-2)-D-O}{(n-1)(n-2)}.
\]
Consequently the target $\overline\delta\le2$ is equivalent, for
these planar arrangements, to
\[
 D+O\ge2(n-2).
\]
This exposes the information erased by the crude double counting:
edges not counted twice and the half-unit diameter savings from odd
polygons. It is a checkable equivalent condition, not a proof of its
lower bound. The cited planar theorem provides the stronger established
answer by using additional arrangement structure.

\paragraph{3. A boundary case and the failed extension.}
In dimension one, the hyperplanes are $n$ distinct points on a line.
The $n-1$ bounded cells are intervals, each with a two-vertex,
one-edge graph; their average diameter is exactly one. This checks
the dimension and cell-count conventions.

In higher dimensions it is invalid to replace a cell graph by a
cycle or identify diameter with half its number of facets. For
example a $d$-simplex has $d+1$ facets but complete graph on $d+1$
vertices and diameter one. General simple polytopes have more complex
graphs. A count of facet incidences therefore needs an independently
proved relationship with shortest paths; the planar identity supplies
none. This is the first obstruction to extending this calculation.

\paragraph{Verification.}
For three lines there is one triangle, $E_b=3$, $D=3$, $O=1$,
and the exact formula gives $(6-3-1)/2=1$, as required.
The denominator is positive for $n\ge3$, the permitted planar range.
All sums omit unbounded cells, but bounded edges adjacent to them
are retained in $D$ with the correct deficit. No Euclidean distances,
cell volumes, or unsupported general diameter theorem enter the proof.

\subsection{Final status}
UNRESOLVED as an independent proof in arbitrary dimension. The
one-dimensional case, planar asymptotic estimate and exact planar
deficit reformulation are verified. The full planar result is already
known from [S3]; no new solution or global percentage is claimed.

\subsection{Sources}
\begin{itemize}
\item[S1] ULAM.AI and the UnsolvedMath Contributors, supplied problems.json, record 3075 (OPG-605), OpenGarden collection. Catalogue identity and source transcription; CC BY 4.0.
\url{https://huggingface.co/datasets/ulamai/UnsolvedMath}.
\item[S2] Open Problem Garden, Average diameter of a bounded cell of a simple arrangement. Original problem and discussion; the mathematical formulation below is expanded from the supplied source record.
\url{http://www.openproblemgarden.org/op/average_diameter_of_a_bounded_cell_of_a_simple_arrangement}.
\item[S3] A. Deza and F. Xie, \emph{Hyperplane Arrangements with Large Average Diameter}, arXiv:0710.0328, 2007. Original research checked 1 October 2026.
\url{https://arxiv.org/abs/0710.0328}.
\end{itemize}
\end{document}
